\Title{Коллектив кафедры АСВК}{Оформление текстов статей в сборник в \LaTeX, версия 1.0 beta}

\section*{Введение}

Этот файл является одновременно примером и инструкцией по написанию статьи для сборника в \LaTeX. 
Шаблон основан на шаблоне для сборника тезисов ВКР 2017 года и может измениться в дальнейшем, но надо ведь с чего-то начинать.
О работе с \LaTeX можно узнать из книг~[\Refn{book1},\Refn{book2}].

\section{Что нужно для сборки PDF}

В архиве примера есть несколько важных файлов:

\begin{itemize}
 \item 
 \texttt{main\_utf.tex} - основной файл документа. Его нужно передавать команде \texttt{pdflatex} и на его базе будет собираться книга;
 \item 
 \texttt{article\_example.tex} - пример файла статьи;
 \item
 \texttt{example\_image.pdf, example\_image.eps} - вставляемое векторное изображение (PDF для \texttt{pdflatex}, EPS для \texttt{latex}).
\end{itemize}

Используемые модули \LaTeX:

\begin{itemize}
 \item indentfirst, url, graphicx, amssymb, amsmath, amscd, bm
 \item inputenc, babel, titlesec
\end{itemize}

\section{Как писать текст}

Текст статьи разбивается на разделы с помощью команды \texttt{\textbackslash section\{Название раздела\}}. Для оформления разделов без номера (например, для введения), ставим звёздочку после слова section: \texttt{\textbackslash section*\{Название раздела\}}.

Символ двойной кавычки <<\verb.".>> является активным, так что если вам очень нужна двойная кавычка, используйте команду \verb.\dq., но лучше пользуйтесь командами \verb.<<. и \verb.>>. для воспроизведения <<типографских>> кавычек.
Знак номера набирайте командой \verb.\No.; также будьте внимательными к трем видам <<черточек>> (дефис, <<среднее>> и <<длинное>> тире).
Кодировки \texttt{Windows-1251} и \texttt{UTF-8} позволяют некоторые из этих
символов набрать без всяких команд; \emph{не делайте этого!}

Формулы можно использовать как обычные: \( (a+b)^2 = a^2 + 2ab + b^2 \), так и выключные:
\[
    \forall \varepsilon>0 \quad
    \exists \delta(\varepsilon)>0
    \; \colon \forall x \; 0<|x-a|<\delta
    \Rightarrow |f(x)-b|<\varepsilon,
\]
но применять при этом автоматическую нумерацию формул нельзя, как и все остальные виды автоматической нумерации: помните, вы делаете фрагмент документа \LaTeX{}, а не документ целиком.

Для того, чтобы продолжать пользоваться ссылками и не портить целостность документа, используйте команды \texttt{\textbackslash Refn\{\}}, \texttt{\textbackslash Pageref\{\}} и \texttt{\textbackslash Label\{\}} вместо стандартных \texttt{\textbackslash ref\{\}}, \texttt{\textbackslash pageref\{\}} и \texttt{\textbackslash label\{\}} (то есть, с большой буквы!). Для удобства работы в современных \LaTeX-редакторах, которые умеют делать автоподстановку имён ссылок, можно писать весь документ со стандартными командами и перед отправкой сделать замену их имён.

Список литературы составляется вручную так, как написано в примере. Обратите внимание на использование команд \texttt{\textbackslash Label\{\}}.

\section{Как вставлять изображения}

Изображения должны быть чёрно-белыми, поскольку дальше это печатается в типографии без цветных вставок.

Крайне желательно использование векторных изображений. Для создания векторных изображений можно использовать специализированные редакторы вроде Inkscape~[\Refn{inkscape}] (доступен для Windows, Linux, Mac OS X). Графики можно строить родными средствами \LaTeX или с помощью других сред вроде gnuplot (важно, чтобы результат был в векторном формате).

Если вы используете gnuplot для построения графиков, рекомендуем использовать эут первую строку:

\texttt{set terminal postscript enhanced eps monochrome solid}

Так сгенерируется файл сразу в нужном формате (EPS чёрно-белый).

Если без растрового изображения никак (фотография, скриншот и т.п.), то нужно использовать форматы PNG, JPEG или TIFF (GIF не поддерживается). Будьте готовы к тому, что изображение будет масштабироваться и может исказиться.

Для включения изображений используется функционал пакета graphicx. Расширение в имени файла не пишется! Нельзя использовать плавающие блоки (\texttt{figure} и \texttt{table}), для таблиц используем \texttt{tabular}, для изображений формируем подписи вручную.

Пример включения изображения:

\texttt{\textbackslash includegraphics[witdh=0.9\textbackslash textwidth]{example\_image}}

\medskip
\includegraphics[width=0.9\textwidth]{example_image}
\begin{centering}
 \medskip
 Рисунок 1. Описание изображения\par
 \medskip
\end{centering}

Изображение должно быть предоставлено в двух форматах: EPS (и для векторных, и для растровых изображений) для обработки с помощью \texttt{latex} и в родном формате: PDF (для векторных изображений), JPEG/PNG/TIFF (для растровых).

Для того, чтобы сконвертировать EPS в PDF в linux, можно использовать утилиту \texttt{eps2pdf}~[\Refn{eps2pdf}]. Для растровых изображений это можно делать с помощьб утилиты \texttt{convert} из \texttt{imagemagick}. Также Inkscape умеет экспортировать изображения в формате EPS.


Для 
\textit{TODO: здесь нужен стандарт для именования файлов, чтобы не перепуталась вся графика между статьями}.

\References
\begin{enumerate}
\item
    \Label{book1}
    Львовский~С.~М.
    \emph{Набор и верстка в \LaTeX.}
    М.:~МЦНМО, 2003.

\item
    \Label{book2}
    Гуссенс~М., Миттельбах~Ф., Самарин А.
    \emph{Путеводитель по пакету \LaTeX{} и его расширению \LaTeX2e.}
    М.: Мир, 1999.
    
\item
    \Label{inkscape}
    Inkscape 
    \texttt{https://inkscape.org/ru/}
    
\item
    \Label{eps2pdf}
    eps2pdf for Linux
    \texttt{https://linux.die.net/man/1/epstopdf}
\end{enumerate}
